\section{Roots of Unity}

Let $r$ denote an $n$th root of unity. If there is no positive integer $m < n$ such that $r^m = 1$, then we say that $r$ is a primitive root $n$th root of unity.

\begin{theorem}
Let $R = \cos(2\pi/n) + i \sin(2\pi/n)$. The number $R^k ,~(1 \le k \le n)$ is a primitive $n$th roof of unity, if and only if, the integers $k$ and $n$ are relatively prime. 
\end{theorem}

\begin{proof}
First assume that $k$ and $n$ have a common factor $d$. Then $n = sd,  k = td$. Then 
\[ (R^k)^s = (R^{td})^s = (R^{sd})^t = (R^n)^t = 1.\]
Since $s< n$, and $(R^k)^s = 1$, it follows that $R^k$ is not a primitive root of unity. \\

We next assume that $k$ and $n$ are relatively prime, and show that $(R^k)^n$ cannot have the value $1$ if $m < n$. We use De Moivre's theorem to write
\[ (R^k)^m = \cos{2km\pi\over n} + i \sin{2km\pi\over n}.\]
If the latter quantity is equal to $1$, then $2km\pi/n$ must be a multiple of $2\pi$,  say $(2 km\pi/n) = 2h\pi$. Therefore $km$ must be divisible by $n$. But this is impossible since $k$ has no
factor in common with $n$ and $m < n$ by assumption. Therefore $(R^k)^m \not= 1$ if $m < n$. Therefore $R^k$ is a primitive root of unity if $k$ and $n$ are relatively prime. 
\end{proof}